
\documentclass[12pt, a4paper, openany, twoside]{book}
\usepackage[left=1in,right=1in,bottom=0.5in,top=0.8in]{geometry}
\usepackage{amsfonts}
\usepackage{amsmath, amssymb}
\usepackage{graphicx}
\usepackage[font=small,labelfont=bf]{caption}
\usepackage{epstopdf}
\usepackage[pdfpagelabels,hyperindex]{hyperref}
\usepackage{xcolor}
\usepackage{amsthm}
\usepackage{float}
\usepackage{pgfplots}
\usepackage{listings}
\usepackage{longtable}
\usepackage{mathrsfs}
\usepackage{mathtools}
\usepackage{hyperref}
\usepackage[most]{tcolorbox}
\usepackage{amssymb}
\usepackage{algorithm}
\usepackage{algpseudocode}
\usepackage{tikz-cd}
\usepackage{comment}
\usepackage{mathpartir}

\linespread{1.3} 

\hypersetup{
pdftitle={.},
pdfauthor={Hojune Lee},
}

\newcommand{\overbar}[1]{\mkern 1.5mu\overline{\mkern-1.5mu#1\mkern-1.5mu}\mkern 1.5mu}
\DeclareRobustCommand{\stirling}{\genfrac\{\}{0pt}{}}
\newtheorem{thm}{Theorem}[section]
\newtheorem{cor}[thm]{Corollary}
\newtheorem{prop}[thm]{Proposition}
\newtheorem{lem}[thm]{Lemma}
\newtheorem{conj}[thm]{Conjecture}
\newtheorem{quest}[thm]{Question}
\newtheorem{ppty}[thm]{Property}
\newtheorem{ppties}[thm]{Properties}
\newtheorem{axiom}[thm]{Axiom}
\newtheorem{claim}[thm]{Claim}
\newtheorem{prob}[thm]{Problem}


\theoremstyle{definition}
\newtheorem{defn}[thm]{Definition}
\newtheorem{defns}[thm]{Definitions}
\newtheorem{con}[thm]{Construction}
\newtheorem{exmp}[thm]{Example}
\newtheorem{exmps}[thm]{Examples}
\newtheorem{notn}[thm]{Notation}
\newtheorem{notns}[thm]{Notations}
\newtheorem{addm}[thm]{Addendum}
\newtheorem{exer}[thm]{Exercise}
\newtheorem{limit}[thm]{Limitation}


\theoremstyle{remark}
\newtheorem{rem}[thm]{Remark}
\newtheorem{rems}[thm]{Remarks}
\newtheorem{warn}[thm]{Warning}
\newtheorem{sch}[thm]{Scholium}
\newenvironment{thmbox}[1][]{\begin{tcolorbox}[colback=yellow!10!white,colframe=red!75!black,title={#1}]}{\end{tcolorbox}}


% newcommand bb
    \newcommand{\BA}{{\mathbb {A}}} \newcommand{\BB}{{\mathbb {B}}}
    \newcommand{\BC}{{\mathbb {C}}} \newcommand{\BD}{{\mathbb {D}}}
    \newcommand{\BE}{{\mathbb {E}}} \newcommand{\BF}{{\mathbb {F}}}
    \newcommand{\BG}{{\mathbb {G}}} \newcommand{\BH}{{\mathbb {H}}}
    \newcommand{\BI}{{\mathbb {I}}} \newcommand{\BJ}{{\mathbb {J}}}
    \newcommand{\BK}{{\mathbb {U}}} \newcommand{\BL}{{\mathbb {L}}}
    \newcommand{\BM}{{\mathbb {M}}} \newcommand{\BN}{{\mathbb {N}}}
    \newcommand{\BO}{{\mathbb {O}}} \newcommand{\BP}{{\mathbb {P}}}
    \newcommand{\BQ}{{\mathbb {Q}}} \newcommand{\BR}{{\mathbb {R}}}
    \newcommand{\BS}{{\mathbb {S}}} \newcommand{\BT}{{\mathbb {T}}}
    \newcommand{\BU}{{\mathbb {U}}} \newcommand{\BV}{{\mathbb {V}}}
    \newcommand{\BW}{{\mathbb {W}}} \newcommand{\BX}{{\mathbb {X}}}
    \newcommand{\BY}{{\mathbb {Y}}} \newcommand{\BZ}{{\mathbb {Z}}}

% newcommand  scr
    \newcommand{\sA}{{\mathscr {A}}} \newcommand{\sB}{{\mathscr {B}}}
    \newcommand{\sC}{{\mathscr {C}}} \newcommand{\sD}{{\mathscr {D}}}
    \newcommand{\sE}{{\mathscr {E}}} \newcommand{\sF}{{\mathscr {F}}}
    \newcommand{\sG}{{\mathscr {G}}} \newcommand{\sH}{{\mathscr {H}}}
    \newcommand{\sI}{{\mathscr {I}}} \newcommand{\sJ}{{\mathscr {J}}}
    \newcommand{\sK}{{\mathscr {K}}} \newcommand{\sL}{{\mathscr {L}}}
    \newcommand{\sN}{{\mathscr {N}}} \newcommand{\sM}{{\mathscr {M}}}
    \newcommand{\sO}{{\mathscr {O}}} \newcommand{\sP}{{\mathscr {P}}}
    \newcommand{\sQ}{{\mathscr {Q}}} \newcommand{\sR}{{\mathscr {R}}}
    \newcommand{\sS}{{\mathscr {S}}} \newcommand{\sT}{{\mathscr {T}}}
    \newcommand{\sU}{{\mathscr {U}}} \newcommand{\sV}{{\mathscr {V}}}
    \newcommand{\sW}{{\mathscr {W}}} \newcommand{\sX}{{\mathscr {X}}}
    \newcommand{\sY}{{\mathscr {Y}}} \newcommand{\sZ}{{\mathscr {Z}}}


% newcommand cal
    \newcommand{\CA}{{\mathcal {A}}} \newcommand{\CB}{{\mathcal {B}}}
    \newcommand{\CC}{{\mathcal {C}}} \newcommand{\CD}{{\mathcal {D}}}
    \newcommand{\CE}{{\mathcal {E}}} \newcommand{\CF}{{\mathcal {F}}}
    \newcommand{\CG}{{\mathcal {G}}} \newcommand{\CH}{{\mathcal {H}}}
    \newcommand{\CI}{{\mathcal {I}}} \newcommand{\CJ}{{\mathcal {J}}}
    \newcommand{\CK}{{\mathcal {K}}} \newcommand{\CL}{{\mathcal {L}}}
    \newcommand{\CM}{{\mathcal {M}}} \newcommand{\CN}{{\mathcal {N}}}
    \newcommand{\CO}{{\mathcal {O}}} \newcommand{\CP}{{\mathcal {P}}}
    \newcommand{\CQ}{{\mathcal {Q}}} \newcommand{\CR}{{\mathcal {R}}}
    \newcommand{\CS}{{\mathcal {S}}} \newcommand{\CT}{{\mathcal {T}}}
    \newcommand{\CU}{{\mathcal {U}}} \newcommand{\CV}{{\mathcal {V}}}
    \newcommand{\CW}{{\mathcal {W}}} \newcommand{\CX}{{\mathcal {X}}}
    \newcommand{\CY}{{\mathcal {Y}}} \newcommand{\CZ}{{\mathcal {Z}}}

    % newcommand frak
     \newcommand{\fa}{{\mathfrak{a}}}  \newcommand{\fb}{{\mathfrak{b}}}
     \newcommand{\fc}{{\mathfrak{c}}}  \newcommand{\fd}{{\mathfrak{d}}}
     \newcommand{\fe}{{\mathfrak{e}}}  \newcommand{\ff}{{\mathfrak{f}}}
     \newcommand{\fg}{{\mathfrak{g}}}  \newcommand{\fh}{{\mathfrak{h}}}
     \newcommand{\fii}{{\mathfrak{i}}}  \newcommand{\fj}{{\mathfrak{j}}}
     \newcommand{\fk}{{\mathfrak{m}}}  \newcommand{\fl}{{\mathfrak{l}}}
     \newcommand{\fm}{{\mathfrak{m}}}  \newcommand{\fn}{{\mathfrak{n}}}
     \newcommand{\fo}{{\mathfrak{o}}}  \newcommand{\fp}{{\mathfrak{p}}}
     \newcommand{\fq}{{\mathfrak{q}}}  \newcommand{\fr}{{\mathfrak{r}}}
     \newcommand{\fs}{{\mathfrak{s}}}  \newcommand{\ft}{{\mathfrak{t}}}
     \newcommand{\fu}{{\mathfrak{u}}}  \newcommand{\fv}{{\mathfrak{v}}}
     \newcommand{\fw}{{\mathfrak{w}}}  \newcommand{\fx}{{\mathfrak{x}}}
     \newcommand{\fy}{{\mathfrak{y}}}  \newcommand{\fz}{{\mathfrak{z}}}

    \newcommand{\fA}{{\mathfrak{A}}}  \newcommand{\fB}{{\mathfrak{B}}}
     \newcommand{\fC}{{\mathfrak{C}}}  \newcommand{\fD}{{\mathfrak{D}}}
     \newcommand{\fE}{{\mathfrak{E}}}  \newcommand{\fF}{{\mathfrak{F}}}
     \newcommand{\fG}{{\mathfrak{G}}}  \newcommand{\fH}{{\mathfrak{H}}}
     \newcommand{\fI}{{\mathfrak{I}}}  \newcommand{\fJ}{{\mathfrak{J}}}
     \newcommand{\fK}{{\mathfrak{K}}}  \newcommand{\fL}{{\mathfrak{L}}}
     \newcommand{\fM}{{\mathfrak{M}}}  \newcommand{\fN}{{\mathfrak{N}}}
     \newcommand{\fO}{{\mathfrak{O}}}  \newcommand{\fP}{{\mathfrak{P}}}
     \newcommand{\fQ}{{\mathfrak{Q}}}  \newcommand{\fR}{{\mathfrak{R}}}
     \newcommand{\fS}{{\mathfrak{S}}}  \newcommand{\fT}{{\mathfrak{T}}}
     \newcommand{\fU}{{\mathfrak{U}}}  \newcommand{\fV}{{\mathfrak{V}}}
     \newcommand{\fW}{{\mathfrak{W}}}  \newcommand{\fX}{{\mathfrak{X}}}
     \newcommand{\fY}{{\mathfrak{Y}}}  \newcommand{\fZ}{{\mathfrak{Z}}}



 % newcommand :rm
     \newcommand{\RA}{{\mathrm {A}}} \newcommand{\RB}{{\mathrm {B}}}
    \newcommand{\RC}{{\mathrm {C}}} \newcommand{\RD}{{\mathrm {D}}}
    \newcommand{\RE}{{\mathrm {E}}} \newcommand{\RF}{{\mathrm {F}}}
    \newcommand{\RG}{{\mathrm {G}}} \newcommand{\RH}{{\mathrm {H}}}
    \newcommand{\RI}{{\mathrm {I}}} \newcommand{\RJ}{{\mathrm {J}}}
    \newcommand{\RK}{{\mathrm {K}}} \newcommand{\RL}{{\mathrm {L}}}
    \newcommand{\RM}{{\mathrm {M}}} \newcommand{\RN}{{\mathrm {N}}}
    \newcommand{\RO}{{\mathrm {O}}} \newcommand{\RP}{{\mathrm {P}}}
    \newcommand{\RQ}{{\mathrm {Q}}} \newcommand{\RR}{{\mathrm {R}}}
    \newcommand{\RS}{{\mathrm {S}}} \newcommand{\RT}{{\mathrm {T}}}
    \newcommand{\RU}{{\mathrm {U}}} \newcommand{\RV}{{\mathrm {V}}}
    \newcommand{\RW}{{\mathrm {W}}} \newcommand{\RX}{{\mathrm {X}}}
    \newcommand{\RY}{{\mathrm {Y}}} \newcommand{\RZ}{{\mathrm {Z}}}

    \newcommand{\Ad}{{\mathrm{Ad}}} \newcommand{\Aut}{{\mathrm{Aut}}}
    \newcommand{\Br}{{\mathrm{Br}}} \newcommand{\Ch}{{\mathrm{Ch}}}
    \newcommand{\cod}{{\mathrm{cod}}} \newcommand{\cont}{{\mathrm{cont}}}
    \newcommand{\cl}{{\mathrm{cl}}}   \newcommand{\Cl}{{\mathrm{Cl}}}
    \newcommand{\disc}{{\mathrm{disc}}}\newcommand{\Eis}{{\mathrm{Eis}}}
    \newcommand{\Div}{{\mathrm{Div}}} \renewcommand{\div}{{\mathrm{div}}}
    \newcommand{\End}{{\mathrm{End}}} \newcommand{\Frob}{{\mathrm{Frob}}}
    \newcommand{\Gal}{{\mathrm{Gal}}} \newcommand{\GL}{{\mathrm{GL}}}
    \newcommand{\Hom}{{\mathrm{Hom}}} \renewcommand{\Im}{{\mathrm{Im}}}
    \newcommand{\Ind}{{\mathrm{Ind}}} \newcommand{\ind}{{\mathrm{ind}}}
    \newcommand{\inv}{{\mathrm{inv}}}
    \newcommand{\Isom}{{\mathrm{Isom}}} \newcommand{\Jac}{{\mathrm{Jac}}}
    \newcommand{\ad}{{\mathrm{ad}}}  \newcommand{\Tr}{{\mathrm{Tr}}}
    \newcommand{\Ker}{{\mathrm{Ker}}} \newcommand{\Ros}{{\mathrm{Ros}}}
    \newcommand{\Lie}{{\mathrm{Lie}}} \newcommand{\Hol}{{\mathrm{Hol}}}

    \newcommand{\cyc}{{\mathrm{cyc}}}\newcommand{\id}{{\mathrm{id}}}
    \newcommand{\new}{{\mathrm{new}}} \newcommand{\NS}{{\mathrm{NS}}}
    \newcommand{\ord}{{\mathrm{ord}}} \newcommand{\rank}{{\mathrm{rank}}}
    \newcommand{\PGL}{{\mathrm{PGL}}} \newcommand{\Pic}{\mathrm{Pic}}
    \newcommand{\cond}{\mathrm{cond}} \newcommand{\Is}{{\mathrm{Is}}}
    \renewcommand{\Re}{{\mathrm{Re}}} \newcommand{\reg}{{\mathrm{reg}}}
    \newcommand{\Res}{{\mathrm{Res}}} \newcommand{\Sel}{{\mathrm{Sel}}}
    \newcommand{\RTr}{{\mathrm{Tr}}} \newcommand{\alg}{{\mathrm{alg}}}
    \newcommand{\PSL}{{\mathrm{PSL}}}

\newcommand{\coker}{{\mathrm{coker}}}
\newcommand{\val}{{\mathrm{val}}} \newcommand{\sign}{{\mathrm{sign}}}
\newcommand{\mult}{{\mathrm{mult}}} \newcommand{\Vol}{{\mathrm{Vol}}}
\newcommand{\Meas}{{\mathrm{Meas}}}\renewcommand{\mod}{\ \mathrm{mod}\ }
\newcommand{\Ann}{\mathrm{Ann}}
\newcommand{\Tor}{\mathrm{Tor}}
\newcommand{\Supp}{\mathrm{Supp}}\newcommand{\supp}{\mathrm{supp}}
\newcommand{\Max}{\mathrm{Max}}
\newcommand{\Coker}{\mathrm{Coker}}
\newcommand{\Stab}{\mathrm{Stab}}
\newcommand{\Irr}{\mathrm{Irr}}\newcommand{\Inf}{\mathrm{Inf}}\newcommand{\Sup}{\mathrm{Sup}}
\newcommand{\rk}{\mathrm{rk}}\newcommand{\Fil}{\mathrm{Fil}}
\newcommand{\Sim}{{\mathrm{Sim}}} \newcommand{\SL}{{\mathrm{SL}}}
\newcommand{\Spec}{{\mathrm{Spec}}} \newcommand{\SO}{{\mathrm{SO}}}
\newcommand{\SU}{{\mathrm{SU}}} \newcommand{\Sym}{{\mathrm{Sym}}}
\newcommand{\sgn}{{\mathrm{sgn}}} \newcommand{\tr}{{\mathrm{tr}}}
\newcommand{\tor}{{\mathrm{tor}}}  \newcommand{\ur}{{\mathrm{ur}}}
\newcommand{\vol}{{\mathrm{vol}}}  \newcommand{\ab}{{\mathrm{ab}}}
\newcommand{\Sh}{{\mathrm{Sh}}} \newcommand{\Ell}{{\mathrm{Ell}}}
\newcommand{\Char}{{\mathrm{Char}}}\newcommand{\Tate}{{\mathrm{Tate}}}
\newcommand{\corank}{{\mathrm{corank}}} \newcommand{\Cond}{{\mathrm{Cond}}}
\newcommand{\Inn}{{\mathrm{Inn}}} \newcommand{\Spf}{{\mathrm{Spf}}}
\newcommand{\Mat}{{\mathrm{Mat}}}


    \font\cyr=wncyr10  \newcommand{\Sha}{\hbox{\cyr X}}
    \newcommand{\wt}{\widetilde} \newcommand{\wh}{\widehat} \newcommand{\ck}{\check}
    \newcommand{\pp}{\frac{\partial\bar\partial}{\pi i}}
    \newcommand{\pair}[1]{\langle {#1} \rangle}
    \newcommand{\wpair}[1]{\left\{{#1}\right\}}
    \newcommand{\intn}[1]{\left( {#1} \right)}
    \newcommand{\norm}[1]{\|{#1}\|}
    \newcommand{\sfrac}[2]{\left( \frac {#1}{#2}\right)}
    \newcommand{\ds}{\displaystyle}
    \newcommand{\ov}{\overline}
    \newcommand{\Gros}{Gr\"{o}ssencharaktere}
    \newcommand{\incl}{\hookrightarrow}
    \newcommand{\lra}{\longrightarrow}
     \newcommand{\ra}{\rightarrow}
    \newcommand{\imp}{\Longrightarrow}
    \newcommand{\lto}{\longmapsto}
    \newcommand{\bs}{\backslash}
    \newcommand{\nequiv}{\equiv\hspace{-7.8pt}/}
    \theoremstyle{plain}


\definecolor{energy}{RGB}{114,0,172}
\definecolor{freq}{RGB}{45,177,93}
\definecolor{spin}{RGB}{251,0,29}
\definecolor{signal}{RGB}{203,23,206}
\definecolor{circle}{RGB}{217,86,16}
\definecolor{average}{RGB}{203,23,206}
\newcommand{\K}{\operatornamewithlimits{K}}
\colorlet{shadecolor}{gray!20}
\pgfplotsset{compat=1.9}
\def\N{10}
\def\M{4}
\usepgflibrary{fpu}


\def\upint{\mathchoice%
    {\mkern13mu\overline{\vphantom{\intop}\mkern7mu}\mkern-20mu}%
    {\mkern7mu\overline{\vphantom{\intop}\mkern7mu}\mkern-14mu}%
    {\mkern7mu\overline{\vphantom{\intop}\mkern7mu}\mkern-14mu}%
    {\mkern7mu\overline{\vphantom{\intop}\mkern7mu}\mkern-14mu}%
  \int}
\def\lowint{\mkern3mu\underline{\vphantom{\intop}\mkern7mu}\mkern-10mu\int}



\makeatletter
\let\c@equation\c@thm
\raggedbottom
\makeatother
\numberwithin{equation}{section}
%--------Meta Data: Fill in your info------
\author{Hojune Lee, 20210541}

\title{Abstract Algebra}
\begin{document}

\begin{titlepage}
    \begin{center}
        \vspace*{9cm}
            
        \Huge
        \textbf{Abstract Algebra}
    
        \vspace{1cm}
        \large
        Summary note of Dummit \& Foote's Abstract Algebra
        \vspace{3cm}
        
        \LARGE
        \textbf{Hojune Lee}
            
        \vspace{8cm}
            
        \normalsize
        \textbf{School of Computing, KAIST}\\  
    \end{center}
\end{titlepage}

\hypersetup{linkcolor=black}
\tableofcontents
\hypersetup{linkcolor=blue}

\newpage 

\chapter{Introduction to Groups}

\section{Groups}

\begin{tcolorbox}[colback=yellow!10!white,colframe=brown!75!black,title=Binary operation]

\begin{enumerate}
    \item A \textit{binary operation} $\star$ on a set $G$ is a map $\star : G \times G \rightarrow G$. We denote $\forall a, b \in G, \star(a, b)$ by $a \star b$. 
    \item A binary operation $\star$ on a set $G$ is \textit{associative} if \[\forall a, b, c \in G, a \star (b \star c) = (a \star b) \star c\]
    \item A binary operation $\star$ on a set $G$ is \textit{commutative} if \[\forall a, b \in G, a \star b = b \star a\]
    \item For a set $G$, $H \subseteq G$ and a binary operation $\star$ on $G$, if the restriction of $\star$ to $H$ is a binary operation on $H$, $H$ is said to be \textit{closed} under $\star$. 
\end{enumerate}

\end{tcolorbox}

\begin{tcolorbox}[colback=yellow!10!white,colframe=brown!75!black,title=Group]

\begin{enumerate}
    \item A \textit{group} is an ordered pair $(G, \star)$ where $G$ is a set and $\star$ is a binary operation on $G$, which satisfies 
    \begin{enumerate}
        \item $\star$ is associative.
        \item $\exists~ e \in G, \forall a \in G, a = a \star e = e \star a$
        \item For $e$ of (b), $\forall a \in G, \exists a^{-1} \in G, a \star a^{-1} = a^{-1} \star a = e$
    \end{enumerate}
    \item Group $(G, \star)$ is called \textit{abelian} when $\star$ is also commutative. 
\end{enumerate}

\end{tcolorbox}

\begin{tcolorbox}[colback=yellow!10!white,colframe=brown!75!black,title=Direct product,breakable]
Let $(A, \star)$ and $(B, \diamond)$ are groups. We can form group $A \times B$ called \textit{direct product} of $A$ and $B$ by
\[A \times B = \{(a, b) | a \in A, b \in B\}\]
with following binary operation $\cdot$ : 
\[(a_1, b_1) \cdot (a_2, b_2) = (a_1 \star a_2, b_1 \diamond b_2)\]

Here, one can claim that $(A \times B, \cdot)$ is a group. The associativity of $\cdot$ is clear. 
And let's denote the identity of $(A, \star), (B, \diamond)$ by $e_1, e_2$ respectively. We can easily check that $(e_1, e_2)$ is identity of $(A\times B, \cdot)$
and $(e_1, e_2) \in A \times B$. This holds group axiom (2). Finally, $\forall a \in A, b \in B$, we can compose $(a^{-1}, b^{-1})$. Then by the definition of $\cdot$, 
$(a, b) \cdot (a^{-1}, b^{-1}) = (a^{-1}, b^{-1}) = (e_1, e_2)$ holds. So, $(a, b)^{-1} = (a^{-1}, b^{-1})$, holds group axiom (3). 
\end{tcolorbox}

\begin{tcolorbox}[colback=yellow!10!white,colframe=red!75!black,title=Propositions on Groups, breakable]
Let $(G, \star)$ be any group. Following always holds. 
\begin{enumerate}
    \item The identity of $G$ is unique.
    \item $\forall a \in G$, $a^{-1}$ is uniquely determined. 
    \item $\forall a \in G, (a^{-1})^{-1} = a$
    \item $\forall a, b \in G, (a \star b)^{-1} = b^{-1} \star a^{-1}$
    \item $\forall a_1, \cdots, a_n \in G$ the value of $a_1 \star a_2 \star \cdots \star a_n$ is independent of how the expression is bracketed.
\end{enumerate}

All the proofs are easy. 

\end{tcolorbox}

\begin{tcolorbox}[colback=yellow!10!white,colframe=red!75!black,title=Cancellation Laws hold in $G$]
\begin{enumerate}
    \item $au = av \implies u = v$
    \item $ub = vb \implies u = v$
\end{enumerate}

\end{tcolorbox}

\begin{tcolorbox}[colback=yellow!10!white,colframe=brown!75!black,title=Group order]

For any group $G$ and $x \in G$, let's say the \textit{order} of $x$ (denote by $|x|$) by 
\[|x| = \underset{n \in \BN}{\text{argmin}}~\{ x^{n} = 1\}\]
If such $n$ does not exists, $|x| = \infty$. 

\end{tcolorbox}

\newpage 

\section{Dihedral Groups}

The definition of \textit{dihedral group} is quite abstracted. 

\begin{tcolorbox}[colback=yellow!10!white,colframe=brown!75!black,title=Dihedral Groups]
One natural, and important structure is symmetries of geometric objects. For $n \in \BZ^{+}, n \geq 3$, the \textit{dihedral group} $D_{2n}$ 
be the set of symmetries on regular $n$-gon, where a symmetry is any rigid motion. 
\end{tcolorbox}

Now let's imagine that we label each vertices of such $n$-gon as $\{1, 2, \cdots, n\}$. 
Then, each symmetries can be viewed a permutation $\sigma$ on $\{1, 2, \cdots, n\}$. 
Then now we can formulate the group $D_{2n}$ with composition $\circ$. Let $s, t \in D_{2n}$ be two symmetries. 
Composition $s\circ t = st$, is defined by apply $t$ and then $s$. 

One interesting thing is that, $|D_{2n}| = 2n$. It can be proved. In every symmetry, vertex 1 and 2 are adjacent. 
So, our job is clear. Choose the position of vertex 1. ($n$ cases) And then, choose the position of vertex 2 which adjacent with vertex 1. (2 cases) So, we have $2n$ cases total, implies that $|D_{2n}| = 2n$. 
Or, we can enumerate what such $2n$ symmetries are. $n$ symmetries are came from rotation by $2\pi i / n, i \in \{0, \cdots, n-1\}$. Remaining $n$ symmetries are came from flip(reflection), with $n$ axis of $n$-gon. 

From now on, we'll fix our setting to understand $D_{2n}$ further. Fix a regular $n$-gon centered at the origin on $\BR^2$ and label the vertices in clockwise manner, from $1$ to $n$. 
And denote $r \in D_{2n}$ be a $\frac{2\pi}{n}$ rotation. And denote $s \in D_{2n}$ be a reflection via axis containing vertex 1. We can observe following propositions. 

\begin{tcolorbox}[colback=yellow!10!white,colframe=green!75!black,title=Propositions on $D_{2n}$]
\begin{enumerate}
    \item $1, r, r^2, \cdots, r^{n-1}$ are distinct, and $|r| = n$
    \item $|s| = 2$
    \item $s \neq r^{i}$ for any $i \in \{0, 1, \cdots, n-1\}$
    \item $sr^i \neq sr^j$ for all $0 \leq i, j \leq n-1$ with $i \neq j$. Consequence of this is that, 
    \[D_{2n} = \{1, r, r^2, \cdots, r^{n-1}, s, sr, sr^2, \cdots, sr^{n-1}\}\]
    \item $rs = sr^{-1}$, So $D_{2n}$ is \textit{non-abelian}
    \item $r^i s = sr^{-i}, 0 \leq i \leq n$
\end{enumerate}

Everythings are not hard to prove. 

\end{tcolorbox}

Now, let's discussion about \textit{generators} and \textit{relations}. As we can see in above study, 
to represent all elements of $D_{2n}$, two element $r, s$ are sufficient. Similarly, we introduce the notions of \textit{generators} and \textit{relations} for arbitrary groups. 
Although we recall both notions later, let's bring the notions of this basis-like things. 

\begin{tcolorbox}[colback=yellow!10!white,colframe=brown!75!black,title=Generator of $G$]

For a group $G$, $S \subset G$ is called a set of \textit{generators} if every element of $G$ can be written 
as a (finite) product of elements of $S$ and their inverses. Also we write $G = \langle S \rangle$ and say that 
$G$ is \textit{generated by} $S$. 

\end{tcolorbox}

For clear understanding, $\BZ = \langle 1 \rangle, D_{2n} = \langle r, s \rangle$. 

\begin{tcolorbox}[colback=yellow!10!white,colframe=brown!75!black,title=Relation in $G$]

Any equations in group $G$ that the generators of $G$ satisfy are called \textit{relations} in $G$.

\end{tcolorbox}

For example, we have 3 relations in $D_{2n}$ : 
\[r^n = 1, s^2 = 1, rs = sr^{-1}\]
And any other relation between elements of $D_{2n}$ can be derived from these three. (later)
By the way, this notion enable us to represent group in simple way. (Only few elements and equations)
This leads us to following definition. 

\begin{tcolorbox}[colback=yellow!10!white,colframe=brown!75!black,title=Presentation of $G$]

Let's imagine that a group $G$ is generated by $S \subseteq G$ and there are some collection of relations $R_1, \cdots, R_m$ such that 
any relation among the $S$ is consequence of them. In this case, we write 
\[G = \langle S | R_1, R_2, \cdots, R_m \rangle\]
and call this by \textit{presentation} of $G$. 

\end{tcolorbox}

In this step, this notion is quite verbose. So, we'll discuss it later of this book. 

\newpage 

\section{Symmetric Groups}

Now let's move on to next step. Actually our goal is abstraction of natural notions. In this section, we'll 
derive set of all bijections(or permutations). 

\begin{tcolorbox}[colback=yellow!10!white,colframe=brown!75!black,title=Symmetric Group on $\Omega$]

Let $\Omega$ be any set. Then define $S_\Omega$ by the set of all bijections from $\Omega$ to $\Omega$. 
We can also understand $S_\Omega$ is set of all permutations of $\Omega$. For the binary operation $\circ$ which is function composition, 
$(S_\Omega, \circ)$ is a group. (Check!) We call $(S_\Omega, \circ)$ by \textit{symmetric group} on $\Omega$. 

\end{tcolorbox}

One natural consideration is that, when $\Omega = \{1, 2, \cdots, n\}$. In this case, we denote the symmetric group by $(S_n, \circ)$ 
and call it \textit{symmetric group of degree $n$}. We already know that $|S_n| = n!$ 

Let's move on. Can we find some interesting observations from each $\sigma \in S_n$? 

\begin{tcolorbox}[colback=yellow!10!white,colframe=brown!75!black,title=Cycle decomposition]

A \textit{cycle} is a string of integers represents the element of $S_n$ which cyclically permutes these integers and fixed all others. For example, a cycle 
$(a_1, \cdots, a_m) \in S_n$ is the permutation such that sends $a_1 \mapsto a_2, \cdots, a_{m-1} \mapsto a_m, a_m \mapsto a_1$ 

In general, we can decompose for each $\sigma \in S_n$ into $k$ cycles, means that 
\[\forall \sigma \in S_n, \sigma = (a_1~a_2\cdots a_{m_1})(a_{m_1 + 1}\cdots a_{m_2})\cdots(a_{m_{k-1} + 1} \cdots a_{m_k})\]

With this, we can easily compute $\sigma(x)$ for any $x \in [n]$. Find the cycle that $x$ is contained, then the immediately right-next element is $\sigma(x)$. 
This decomposition is called \textit{cycle decomposition of $\sigma$}

\end{tcolorbox}

We can easily obtain the cycle decompositioning algorithm, so skip. And with this, we can easily compute cycle decomposition of $\sigma^{-1}$. 
Just write the numbers in each cycle of the cycle decomposition of $\sigma$ in reverse order. i.e., 
\[\sigma^{-1} = (a_{m_1}~ a_{m_1 - 1} \cdots, a_1) (a_{m_2}~\cdots a_{m_1 + 1})\cdots (a_{m_k} \cdots a_{m_{k-1} + 1})\]

Also, we can observe that $\forall n \geq 3, S_n$ is not abelian group. This is because that $(1~3)\circ(1~2) \neq (1~2)\circ (1~3)$. 
And note that, we can view $\sigma$ by composition of each cycles. (since cycles are permutations of $S_n$ itself) 
Moreover, they are independent to order. i.e., \textit{disjoint cycles commute}. And also, each cycles are independent to cyclic permutation of inner elements. i.e., $(1~2~3) = (2~3~1) = (3~1~2)$

\begin{tcolorbox}[colback=yellow!10!white,colframe=violet!75!white,title=Exercise 10,breakable]

Prove that if $\sigma$ is the $m$-cycle $(a_1, a_2, \cdots, a_m)$, then for all $i \in \{1, 2, \cdots m\}, \sigma^i(a_k) = a_{k+i}$
, where $k + i$ replaced by its least positive residue mod $m$. Deduce that $|\sigma| = m$. 

\begin{proof}
    
Let's denote the least positive residue mod $m$ by $R^+_m : \BZ \rightarrow [m]$. Note that, 
$R^+_m(m) = m$, and $R^+_m(m + 1) = 1$. Then we can immediately obtain that $R^+_m(x + m) = R^+_m(x)$. 
Then, let's prove by mathematical induction w.r.t. $i$. For $i = 1$, 
\[\sigma(a_k) = \begin{cases}
    a_{k+1} & 1 \leq k < m \\
    a_1 & k = m 
\end{cases}\]
However, \[R^+_m(k + 1) = \begin{cases}
    k + 1 & 1 \leq k < m \\
    1 & k = m 
\end{cases} \] 
by definition. 
So, $\sigma(a_k) = a_{R^+_m(k+1)}$ so the base case holds. Now, suppose that 
\[\sigma^i(a_k) = a_{R^+_m(k + i)}\]
Then, 
\[\sigma^{i+1}(a_k) = \sigma (\sigma^i(a_k)) = \sigma(a_{R^+_m(k + i)}) = \begin{cases}
    a_{R^+_m(k + i) + 1} & 1 \leq R^+_m(k + i) < m \\
    a_1 & R^+_m(k + i) = m
\end{cases}\]
However, for $1 \leq R^+_m(k + i) < m$, $R^+_m(k + i) + 1 = R^+_m(k + i + 1)$ and 
$R^+_m(k + i) = m$, then $ 1 = R^+_m(k + i + 1)$. So we can write that 
\[\sigma^{i+1}(a_k) = a_{R^+_m(k + i + 1)}\]
This completes the proof.
\end{proof}

\end{tcolorbox}

\begin{tcolorbox}[colback=yellow!10!white,colframe=violet!75!white,title=Exercise 11,breakable]

Let $\sigma$ be the $m$-cycle $(1, 2, \cdots, m)$. Show that $\sigma^i$ is also an $m$-cycle if and only if $i$ is relatively prime to $m$. 

\begin{proof}
    
We can imagin $\sigma$ as $a_k = k$ in Exercise 10. Then here, we can obatain that 
\[\sigma^i(k) = R^+_m(k + i)\]

So that, we can construct cycle decomposition of $\sigma^i$. The first cycle will be: 
\[(1~R^+_m(1 + i)~R^+_m(R^+_m(1 + i) + i)~\cdots )\]
It can be re-written: 
\[(1~R^+_m(1 + i)~R^+_m(1 + 2i)~\cdots~R^+_m(1 + (l-1)i))\]
Here $l$ is the length of this cycle, and $\sigma^i(R^+_m(1 + (l-1)i)) = 1$. 
This implies that $R^+_m(1 + li) = 1$. Note that $1 \leq l \leq m$ and $1 \leq k < l \implies R^+_m(1 + ki) \neq 1$. 
This immediately implies that $l$ is the least positive integer s.t. $1 + li = 1 + nm$. 
Let $d = \gcd(m, i)$. Then, $l = m/d$. So, $l = m \iff d = 1$. This completes the proof.

\end{proof}

\end{tcolorbox}

We can improve above result into nice theorem. 

\begin{tcolorbox}[colback=yellow!10!white,colframe=red!75!black,title=Corollary: i-th power of cycle]

Let $\sigma$ be the $m$-cycle. For $i \in \BZ$, let $d = \gcd(m, i)$. Then, $\sigma^i$ is decomposed into 
$d$ of length $(m/d)$-cycles. 

\end{tcolorbox}

\begin{tcolorbox}[colback=yellow!10!white,colframe=violet!75!white,title=Exercise 13,breakable]

Show that an element has order 2 in $S_n$ if and only if its cycle decomposition is a product of commuting 2-cycles. 

\begin{proof}
    It can be easily obtained, since we already know that 
    disjoint cycles commute, and for $m$-cycle $\sigma$, $|\sigma| = m$. 
\end{proof}

\end{tcolorbox}

\begin{tcolorbox}[colback=yellow!10!white,colframe=violet!75!white,title=Exercise 14]

Let $p$ be a prime. Show that an element has order $p$ in $S_n$ if and only if its cycle decomposition is product of 
commuting $p$-cycles. Show by an explicit example that this need not be the case if $p$ is not prime. 

\begin{proof}
    
($\Rightarrow$) Suppose that there exists $\sigma'$ in the cycle decomposition of $\sigma$ whose order is $p$ such that 
$\sigma' = (a_1~\cdots~a_n)$ and $n \neq p$. Then here, since cycle decomposition is constructed by disjoint cycles, 
\[\sigma^p(a_k) \mapsto a_{R^+_n(k + p)}\]
And since $p$ is prime, this never be $a_k$. So that, $\sigma^p$ is not identity permutation on $a_k$, contradicts that $|\sigma| = p$. 

However, if $p$ need not to be a prime, we can't guarantee that $R^+_n(k + p) \neq k$. It can be possible when $p$ is multiple of $n$. Let's imagine 
\[\sigma = (1~2)(3~4~5~6)\]
This have order 4. This completes the proof.

\end{proof}

\end{tcolorbox}

\begin{tcolorbox}[colback=yellow!10!white,colframe=red!75!black,title=Proposition 1.3.1.,breakable]

The order of $\sigma \in S_n$ equals the least common multiple of the lengths of the cycles in its cycle decomposition.

\begin{proof}
    
Let $\sigma = \sigma_1 \sigma_2 \cdots \sigma_k$ be the cycle decomposition of $\sigma$. Then since $\sigma_i, \sigma_j$ commute, so 
\[\sigma^i = \sigma_1^i \sigma_2^i \cdots \sigma_k^i\]
To this become identity permutation, one is clear that 
\[\sigma_j^i = I \quad \forall 1 \leq j \leq k \]
Since we know $|\sigma_j| = $ length of $\sigma_j$ denote by $l(j)$, 
$i = l(j) \cdot n_j , \quad \forall 1 \leq j \leq k$. 
To $i$ be an order, $i$ must be minimum positive one, so it is just definition of least common multiple of $l(j)$. 

\end{proof}

\end{tcolorbox}

Here, I note that something I realized: 
For $\sigma \in S_m$, we can find the cycle decomposition of $\sigma$. 
There are length $1$ upto $m$ cycles. Let $k_i$ be the number of $i$-cycles. 
Then, following holds: 
\[m = \sum_{i=1}^{p} i \cdot k_i, \quad |\sigma| = \text{lcm}(i | k_i > 0)\]
And if the $\{k_i\}$ fixed, we can find the number of cases by 
\[\frac{p!}{\sum_{i = 1}^{p}i^{k_i} \cdot k_i!}\]

\newpage 

\section{Matrix Groups}

Actually, field is our later interests. However, to understand group theory, very powerful tool is understanding linear algebra. 
So we alert matrix groups, whose coefficients come from fields. For sake of goal, let's describe what the fields are. 

\begin{tcolorbox}[colback=yellow!10!white,colframe=brown!75!black,title=Fields]

A \textit{field} is a set $F$ with two \textbf{commutative} binary operations $+, \cdot$ on $F$ such that 
$(F, +)$ forms an abelian group with identity $0$, and $(F\setminus \{0\}, \cdot)$ also forms an abelian group with identity $1$, 
and following \textit{distributive law} holds: 
\[a \cdot (b + c) = (a \cdot b) + (a \cdot c), \forall a,b,c \in F\]
And we denote $F^{\times} = F \setminus \{0\}$ 

\end{tcolorbox}

Actually our interests are groups of matrices over $F$. Let's define following important group here. 

\begin{tcolorbox}[colback=yellow!10!white,colframe=brown!75!black,title=Matrix groups($GL_n(F)$)]

\[GL_n(F) = \{A | A \in Mat_{n\times n}(F), \det(A) \neq 0, \text{i.e., invertible}\}\]
Let binary operation $\cdot$ denote matrix multiplication. It is easy to see that $(GL_n(F), \cdot)$ is a group. 
It called \textit{general linear group of degree $n$}.  

\end{tcolorbox}

\newpage 

\section{Quaternion Group}

In mathematics, there are some famous counter-examples when we conject some properties on mathematical structures.
This group is sort of them. 

\begin{tcolorbox}[colback=yellow!10!white,colframe=brown!75!black,title=Quaternion Group]

The \textit{quaternion group} $Q_8$ is defined by 
\[Q_8 := \{1, -1, i, -i, j, -j, k, -k\}\]

Here, we define binary operation $\cdot$ as follows; 
\[\forall q \in Q_8, 1 \cdot q = q \cdot 1 = q\]
\[(-1) \cdot (-1) = 1, \quad \forall q \in Q_8, (-1) \cdot q = q \cdot (-1) = -q\]
\[i \cdot i = j \cdot j = k \cdot k = -1\]
\[i \cdot j = k \quad j \cdot i = -k \]
\[j \cdot k = i \quad k \cdot j = -i \]
\[k \cdot i = j \quad i \cdot k = -j\]

\end{tcolorbox}

(I'll add more intuition about this group after I understand well.)

\newpage 

\section{Homomorphisms and Isomorphisms}

By defining groups, we abstracted many mathematical structures into group theories. 
Then now, we can imagine the maps between them and what the \lq look the same group' is.
First, let's define the structure-preserved maps on groups. 

\begin{tcolorbox}[colback=yellow!10!white,colframe=brown!75!black,title=Homomorphism]

Let $(G, \star), (H, \diamond)$ be groups. A map $\varphi : G \rightarrow H$ such that 
\[\forall x, y \in G, \varphi(x \star y) = \varphi(x) \diamond \varphi(y)\]
is called \textit{homomorphism} from $G$ to $H$. 

\end{tcolorbox}

Then now, we define more strong (and more precise to our goal) concepts: Two groups are conceptually equal. 

\begin{tcolorbox}[colback=yellow!10!white,colframe=brown!75!black,title=Isomorphism]

Let the map $\varphi : G \rightarrow H$ be an homomorphism from $G$ to $H$. We call that $\varphi$ is an \textit{isomorphism} when $\varphi$ is a bijection. 
If such isomorphism exists, then we denote $G \cong H$ and say that $G$ and $H$ are \textit{isomorphic}.  

\end{tcolorbox}

On the high level, $G \cong H$ if there exists bijective structure-preserve map $\varphi$. And if $G \cong H$, they are 
exactly same group except that how we write the elements and operation. 

One of the most familar example is that $(\BR, +) \cong (\BR^+, \times)$. Such isomorphism is explicitly given as 
$\varphi(x) = e^x$. It is bijective since $\varphi^{-1}(a) = \log a$. And clearly homomorphism, since $e^{x+y} = e^x e^y$. 
So, both are exactly same structure. 

Following fact is quite not obvious: 

\begin{tcolorbox}[colback=yellow!10!white,colframe=red!75!black,title=First isomorphim classes]
    
Any non-abelian group of order 6 is isomorphic to $S_3$. Further, every group whose order is $6$ 
is isomorphic to one of two groups: $S_3, \BZ / 6\BZ$ and so, $S_3 \ncong \BZ / 6\BZ$. 

\end{tcolorbox}

One of main goal in group theory is that learn skills to prove this kind of statements. 

Now, I add some useful homomorphism criterion here. Suppose that $G$ is generated by 
\(\{s_1, s_2, \cdots, s_m\}\) and also has relations between them. Let $H$ be the another group and 
$r_1, \cdots, r_m \in H$. If all the relations in $G$ is satisfied in $H$ when replace $r_i$ into $s_i$, 
then there is a unique homomorphism $\varphi: G \rightarrow H$ which $\varphi(s_i) = r_i$. Further, if $H$ is 
generated by $\{r_i\}$, then $\varphi$ is surjective. Also, if $|G| = |H| < \infty$, then $\varphi$ is injective. So in this case 
we can also determine $G \cong H$. 

Here, I put some interesting Exercise problems that I proved. 

\begin{tcolorbox}[colback=yellow!10!white,colframe=violet!75!white,title=Exercise Ch 1.6.,breakable]

\begin{enumerate}
    \item If $\varphi : G \rightarrow H$ is an isomorphism, prove that $|\varphi(x)| = |x|$ for all $x \in G$. 
    Deduce that any two isomorphic groups have same number of elements of order $n$ for each $n \in \BZ^+$. 
    Is the result true if $\varphi$ is only assumed to be a homomorphism?
    \item If $\varphi : G \rightarrow H$ is an isomorphism, prove that $G$ is abelian if and only if $H$ is abelian. If $\varphi : G \rightarrow H$ is a homomorphism, what additional conditions are sufficient to ensure that if $G$ is abelian then so is $H$?
    \item Prove that if $|\Delta| = |\Omega|$ then $S_{\Delta} \cong S_\Omega$
    \item $\varphi : G \rightarrow H$ is homomorphism. Then $\varphi(G) < H$ and if $\varphi$ is injective then $G \cong \varphi(G)$. 
    \item $\varphi : G \rightarrow H$ is homomorphism. Define $\ker(\varphi) = \{g \in G | \varphi(g) = 1\}$. Prove that $ker(\varphi) < G$ and $ker(\varphi) = \{1_G\}$ iff $\varphi$ is injective. 
    \item Let $G$ be any group. $\varphi : G \rightarrow G$ define by $g \mapsto g^{-1}$. Prove that $\varphi$ is homomorphism iff $G$ is Abelian.
    \item Same question as above, by $g \mapsto g^2$. 
    \item $A$ is abelian group and fix $k \in \BZ$. $a \mapsto a^k$ is homomorphism from $A$ to $A$. 
    \item Let $G$ be a finite group which possesses an automorphism $\sigma$ such that $\sigma(g) = g$ iff $g = 1$. If $\sigma^2$ is the identity map, then $G$ is abelian. Hint: $g \mapsto x^{-1}\sigma(x)$. 
    \item $G$ is finite group. $x, y \in G, |x| = |y| = 2$ and $G = \langle x, y \rangle$. For $n = |xy|$, $G \cong D_{2n}$. 
\end{enumerate}


\end{tcolorbox}

\newpage 

\section{Group Actions}

Group actions are natural concepts, as we can imagine the $GL_n(\BR)$ as a linear transformation on $\BR^n$ space. 
They \lq act' on vector in $\BR^n$, while preserves the group structures between such linear transformations. So we can define group action very naturally:

\begin{tcolorbox}[colback=yellow!10!white,colframe=brown!75!black,title=Group Action]

A \textit{(left) group action} of a group $G$ on a set $A$ is a map $G \times A \rightarrow A$ written as $g \cdot a$, for every $g \in G, a \in A$, such that 
\begin{enumerate}
    \item $g_1 \cdot (g_2 \cdot a) = (g_1 \cdot g_2) \cdot a, \forall g_1, g_2 \in G, a \in A$
    \item $1 \cdot a = a, \forall a \in A$
\end{enumerate}

\end{tcolorbox}

Now, we'll make quite interesting observation. That's the relations between group actions and permutations. 

\begin{tcolorbox}[colback=yellow!10!white,colframe=green!75!black,title=Group Action as Permutation]

For every $g \in G$, we can define a map 
\[\sigma_g : A \rightarrow A\]
\[\sigma_g : a \mapsto g \cdot a\]
Then, $\sigma_g \in S_A$ and $\varphi_G : g \mapsto \sigma_g$ is a homomorphism of $G \rightarrow S_A$.

\end{tcolorbox}

This $\varphi: G \rightarrow S_A$ is called permutation representation associated to the given action $G$ on $A$.
i.e., If we have $G$ action on $A$, we can find homomorphism from $G$ to $S_A$ always. 
Another direction is also hold. If we have $\varphi : G \rightarrow S_A$ a homomorphism, then 
we can define $G$-action on $A$. 

\begin{tcolorbox}[colback=yellow!10!white,colframe=green!75!black,title=faithful\, kernel\, left regular action]

We say that a $G$ action on $A$ is \textit{faithful} if distinct elements of $G$ induce distinct permutations of $A$. 
We say that the \textit{kernel} of $G$ action on $A$ be $\{g \in G ~|~ g \cdot a = a, ~~ \forall a \in A\}$
Let $A$ be a set $A = G$. Define action of $G$ on $A$ by $g \cdot a = ga$. This action is called \textit{left regular action} of $G$ on itself.

\end{tcolorbox}

\begin{tcolorbox}[colback=yellow!10!white,colframe=violet!75!white,title=Exercise Sec 1.7.,breakable]

\begin{enumerate}
    \item Prove that a group $G$ acts faithfully on a set $A$ if and only if the kernel of the action is $\{1_G\}$. 
    \item Assume $n$ is an even positive integer and show that $D_{2n}$ acts on the set consisting of pairs of opposite vertices of a regular $n$-gon. Find the kernel of this action.
    \item Let $G$ be any group and $A = G$. Show that $g \cdot a = gag^{-1}$ is group action of $G$ on $A$. 
    \item So, each $g \in G$ define a map $x \mapsto gxg^{-1}$. Prove that this is an isomorphism from $G$ to itself. 
    \item Let $H$ be a group acting on a set $A$. Prove that the relation $\sim$ on $A$ defined by 
    \[a \sim b \iff a = hb \text{ for some $h \in H$ }\]
    is an equivalence relation. These equivalence classes is called \textit{orbit} of $x$ under the action $H$.
    \item If $G$ is a finite group and $H < G$ then $|H|$ divides $|G|$. (Lagrange's Theorem)
    \item Group of rigid motions of a tetrahedron is isomorphic to a subgroup of $S_4$.  
\end{enumerate}

\end{tcolorbox}

\newpage 

\chapter{Subgroups}

\section{Definition and Examples}

\begin{tcolorbox}[colback=yellow!10!white,colframe=brown!75!black,title=Subgroup]

Let $G$ be a group. $H \subseteq G$ is called \textit{subgroup} of $G$ if 
$H$ is nonempty and $H$ is closed under products and inverses of $G$. We denote $H \leq G$. 

\end{tcolorbox}

Every subgroup must contain $1_G$. It can be the most easy assertation when $H \leq G$. 

\begin{tcolorbox}[colback=yellow!10!white,colframe=green!75!black,title=Subgroup Criterion]

$H \leq G$ if and only if 
\begin{itemize}
    \item $H \neq \emptyset$, i.e. $1_G \in H$
    \item $\forall x, y \in H, xy^{-1} \in H$
\end{itemize}
If $H$ is finite, then it is sufficient to show that $H$ is nonempty and closed under multiplication.

\end{tcolorbox}

\begin{tcolorbox}[colback=yellow!10!white,colframe=violet!75!white,title=Exercise Sec 2.1.,breakable]

\begin{enumerate}
    \item Give an explicit example of group $G$ and an infinite subset $H$ of $G$ that is closed under the group operation but is not a subgroup of $G$.
    \item Prove that $G$ cannot have a subgroup $H$ with $|H| = n - 1$, where $n = |G| > 2$
    \item Let $G$ be an abelian group. Prove that $\{g \in G ~|~ |g| < \infty\}$ is a subgroup of $G$ (called the \textit{torsion subgroup} of $G$). Give an explicit example where this set is not a subgroup when $G$ is non-abelian.
    \item Let $H$ and $K$ be subgroups of $G$. Prove that $H \cup K$ is a subgroup of $G$ if and only if either $H \subseteq K$ or $K \subseteq H$.
    \item Let $G = GL_n(F)$, where $F$ is any field. Define 
    \[SL_n(F) = \{A \in GL_n(F)~|~\det(A) = 1\}\]
    (called the \textit{special linear group}). Prove that $SL_n(F) \leq GL_n(F)$.
    \item Prove that the intersection of an arbitrary nonempty collection of subgroups of $G$ is again a subgroup of $G$. (Do not assume the collection is countable)
    \item Let $H$ be an subgroup of the additive group of rational numbers with the property that $1/x \in H$ for every nonzero element $x$ of $H$. Prove that $H = 0$ or $\BQ$
    \item Let $H_1 \leq H_2 \leq \cdots $ be can ascending chain of subgroups of $G$. Prove that $\cup_{i=1}^{\infty} H_i$ is a subgroup of $G$.
\end{enumerate}

\end{tcolorbox}

\newpage 

\section{Centralizers, Normalizers, Stabilizers, Kernels}

\begin{tcolorbox}[colback=yellow!10!white,colframe=brown!75!black,title=Centralizer]

Let $G$ be a group and $A \subseteq G$. Define the \textit{centralizer} of $A$ in $G$ by 
\[C_G(A) := \{g \in G ~|~ \forall a \in A, a = gag^{-1}\}\]
In other words, $C_G(A)$ is subset of $G$ such that commutes with every element of $A$.

\end{tcolorbox}

\begin{lem}

For any $A \subseteq G$, $C_G(A) \leq G$

\end{lem}

\begin{proof}
    
Clearly $1_G \in C_G(A)$, so $C_G(A) \neq \emptyset$ for any $A$. Suppose that $x, y \in C_G(A)$. Then, 
\begin{align*}
    \forall a \in A, & (xy^{-1}) a (xy^{-1})^{-1} \\
    & = (xy^{-1})a(yx^{-1}) \\
    &= x (y^{-1}ay) x^{-1} \\
    &= a
\end{align*}
Hence, $C_G(A) \leq G$. 

\end{proof}

\begin{tcolorbox}[colback=yellow!10!white,colframe=brown!75!black,title=Center]

Define the \textit{center} of $G$ by 
\[Z(G) = C_G(G) = \{g \in G ~|~ \forall x \in G, gx = xg \}\]
So we can treat this as a special case of centralizer. 

\end{tcolorbox}

\begin{tcolorbox}[colback=yellow!10!white,colframe=brown!75!black,title=Normalizer]

Let denote $gAg^{-1} = \{gag^{-1}~|~ a \in A\}$. Define the \textit{normalizer} of $A$ in $G$ by 
\[N_G(A) = \{g \in G ~|~ gAg^{-1} = A\}\]
Note that, it does not need to exactly commute with each elements of $A$.

\end{tcolorbox}

\begin{lem}
    For any $A \subseteq G, N_G(A) \leq G$. 
\end{lem}
\begin{proof}
    Same to above. 
\end{proof}

We defined centralizer and normalizer for the \lq subsets' of $G$. We can generalize this into notion of $G$-action on \lq any set'. 

\begin{tcolorbox}[colback=yellow!10!white,colframe=brown!75!black,title=Stabilizer]

Suppose that $G$ is a group acting on set $S$. For fixed $s \in S$, we define \textit{stabilizer} of $s$ in $G$ by 
\[G_s = \{g \in G ~|~ g \cdot s = s\}\]

\end{tcolorbox}

\begin{lem}
Let $G$ be a group acting on set $S$. For $s \in S$, $G_s \leq G$. 
\end{lem}

\begin{proof}
    
The axiom of group action gives us that $1_G \cdot s = s$, so $1_G \in G_s$. And for $g_1, g_2 \in G_s$, 
\[(g_1 g_2^{-1}) \cdot s = g_1 \cdot (g_2^{-1} \cdot s) = g_1 \cdot (g_2^{-1} \cdot (g_2 \cdot s)) = g_1 \cdot s = s\]
So $g_1 g_2^{-1} \in G_s$, hence $G_s \leq G$. 

\end{proof}

\begin{tcolorbox}[colback=yellow!10!white,colframe=brown!75!black,title=Kernel]

Let $G$ be a group acting on set $S$. The \textit{kernel} of left $G$-action on $S$ is defined by 
\[\{g \in G ~|~ \forall s \in S, g \cdot s = s\}\]

\end{tcolorbox}

\begin{tcolorbox}[colback=yellow!10!white,colframe=green!75!black,title=Connections between stabilizer and normalizer]

Suppose that $G$ acts on $\CP(G)$, collection of all subsets of $G$. 
Then action is defined by: 
\[G \times \CP(G) \rightarrow \CP(G)\]
By conjugation action:
\[g: S \mapsto \{gsg^{-1} ~|~ s \in S\}\]

Then stabilizer of $S$ in $G$ is normalizer of $S$ in $G$. 

Similarly, suppose that $N_G(A)$ acts on $A \subseteq G$ by conjugation: 
\[n: a \mapsto nan^{-1}\]
Then, kernel of this action is center of $A$ in $G$. 

\end{tcolorbox}

\begin{tcolorbox}[colback=yellow!10!white,colframe=violet!75!white,title=Exercise Sec 2.2.,breakable]

\begin{enumerate}
    \item Prove that $C_G(Z(G)) = G$ and deduce that $N_G(Z(G)) = G$.
    \item Prove that if $A \subseteq B \subseteq G$ then $C_G(B) \leq C_G(A)$.
    \item Show that if $H \leq G$ then $H \leq N_G(H)$. Give an example that this is not necessarilly true if $H$ is not a subgroup of G. 
    \item Show that if $H \leq G$ then $H \leq C_G(H)$ if and only if $H$ is abelian. 
    \item Let $n \in \BZ$ with $n \geq 3$. $Z(D_{2n}) = 1$ if $n$ is odd, and $Z(D_{2n}) = \{1, r^k\}$ if $n = 2k$.
    \item Let $G = S_n$, fix an $i \in [n]$ and let $G_i = \{\sigma \in G ~|~ \sigma(i) = i\}$ (the stabilizer of $i$ in $G$). Use group actions to prove that $G_i$ is a subgroup of $G$. Find $|G_i|$. 
    \item For any subgroup $H$ of $G$ and any nonempty subset $A$ of $G$ define $N_H(A)$ to be the set $\{h \in H ~|~ hAh^{-1} = A\}$. Show that $N_H(A) = N_G(A) \cap H$ and deduce that $N_H(A)$ is a subgroup of $H$. 
    \item Let $H$ be a subgroup of order 2 in $G$. Show that $N_G(H) = C_G(H)$. Deduce that if $N_G(H) = G$ then $H \leq Z(G)$. 
    \item Prove that $Z(G) \leq N_G(A)$ for any subset $A$ of $G$. 
\end{enumerate}

\end{tcolorbox}

\newpage 

\section{Cyclic groups and Cyclic Subgroups}

\begin{tcolorbox}[colback=yellow!10!white,colframe=brown!75!black,title=Cyclic Group]

A group $H$ is \textit{cyclic} if $H$ can be generated by a single element, i.e., 
\[\exists x \in H, H = \{x^n ~|~ n \in \BZ\}\]
We denote $H = \langle x \rangle$ and $H$ is \textit{generated} by $x$. Cyclic group is always Abelian. 
\end{tcolorbox}

\begin{tcolorbox}[colback=yellow!10!white,colframe=green!75!black,title=Proposition]

If $H = \langle x \rangle$, then $|H| = |x|$. More specifically, 
\begin{itemize}
    \item if $|H| = n < \infty$, then $x^n = 1$ and $1, x, \cdots, x^{n-1}$ are all distinct elements of $H$
    \item if $|H| = \infty$, then $x^n \neq 1$ for all $n \neq 0$ and $x^a \neq x^b$ for all $a \neq b$ in $\BZ$. 
\end{itemize}

\end{tcolorbox}

\begin{proof}
    
Suppose that $|x| = n < \infty$. Then, $\{1, x, \cdots, x^{n-1}\} \subseteq H$ and they are distinct. 
However, for $x^k \in H$, $x^k = x^{np + q} = (x^n)^p \cdot x^q = x^q \in \{1, x, \cdots, x^{n-1}\}$. So, $H = \{1, x, \cdots, x^{n-1}\}$ and $|H| = n$. 
Suppose that $|x| = \infty$. Then, for $a \neq b$, if $x^a = x^b$ then $x^{b - a} = 1$ so order $|x| \leq |b-a| < \infty$ contradiction. So, $|H| = \infty$. 

\end{proof}

\begin{tcolorbox}[colback=yellow!10!white,colframe=green!75!black,title=Proposition]

Let $G$ be an arbitrary group, $x \in G$ and $m, n \in \BZ$. If $x^n = 1, x^m = 1$ then 
$x^d = 1$, where $d = (n, m)$. In particular, if $x^m = 1$ for some $m \in \BZ$ then $|x|$ divides $ m$

\end{tcolorbox}

\begin{proof}
    
Euclidean algorithm gives us that $\exists r, s \in \BZ, d = (n, m) = nr + ms$. Then, 
\[x^d = x^{(n,m)} = x^{nr + ms} = x^{nr} \cdot x^{ms} = 1\]
Suppose that $|x| = n$. Then, $x^m = x^{np + q} = x^q = 1$ for $0 \leq q < n$. Since $|x| = n$, $q = 0$. So, $m = np$, implies that $|x|$ divides $m$. 

\end{proof}

\begin{tcolorbox}[colback=yellow!10!white,colframe=red!75!black,title=Theorem]

Any two cyclic groups of the same order are isomorphic. More specifically, 
\begin{itemize}
    \item if $n \in \BZ^+$ and $\langle x \rangle, \langle y \rangle$ are both cyclic groups of order $n$, then 
    \[\varphi: \langle x \rangle \rightarrow \langle y \rangle\]
    \[x^k \mapsto y^k\]
    is well defined and is an isomorphism. 
    \item if $\langle x \rangle$ is an infinite cyclic group, the map 
    \[\varphi: \BZ \rightarrow \langle x \rangle\]
    \[k \mapsto x^k\]
    is well defined and is an isomorphism.
\end{itemize}

\end{tcolorbox}

\begin{proof}
    
To show that $\varphi$ is well-defined, sufficient to show that 
\[\forall r, s \in \BZ, \quad x^r = x^s \implies \varphi(x^r) = \varphi(x^s)\]
Here, $x^{r-s} = 1$ implies that for $|x| = |y| = n$, $n | (r - s)$. Then, $r = nk + s$ holds. So, 
\[\varphi(x^{r}) = \varphi(x^{nk + s}) = y^{nk + s} = y^s = \varphi(x^s)\]
So $\varphi$ is well-defined. Also, it is easy to show that $\varphi$ is a homomorphism from $\langle x \rangle$ to $\langle y \rangle$. 
Finally, this map is surjective and both groups have same finite order, this is an isomorphism. For the second one, do similarly. 

\end{proof}

From here, we'll denote that $Z_n$ be the order $n$ cyclic group, since all of them are isomorphic. 
Also, we can think that $Z_n \cong \BZ / n\BZ$. However, a given cyclic group may have more than one generator. 

\begin{tcolorbox}[colback=yellow!10!white,colframe=green!75!black,title=Proposition]

Let $G$ be a group, $x \in G$ and let $a \in \BZ - \{0\}$. 

\begin{itemize}
    \item If $|x| = \infty$, then $|x^a| = \infty$. 
    \item If $|x| = n < \infty$, then $|x^a| = \frac{n}{(n, a)}$
    \item In particular, if $|x| = n < \infty$ and $a$ is a positive integer dividing $n$, then $|x^a| = \frac{n}{a}$.  
\end{itemize}

\end{tcolorbox}

\begin{proof}
    
Only second one is meaningful to prove. 

\end{proof}

\begin{tcolorbox}[colback=yellow!10!white,colframe=green!75!black,title=Proposition]

Let $H = \langle x \rangle$. 
\begin{itemize}
    \item $|x| = \infty$ then $H = \langle x^a \rangle$ if and only if $a = \pm 1$
    \item $|x| = n < \infty$ then $H  = \langle x^a \rangle$ if and only if $(a, n) = 1$. 
\end{itemize}

\end{tcolorbox}

\begin{tcolorbox}[colback=yellow!10!white,colframe=red!75!black,title=Theorem]

Let $H = \langle x \rangle$ be a cyclic group. 

\begin{enumerate}
    \item Every subgroup of $H$ is cyclic. More precisely, if $K \leq H$ then either $K = \{1\}$ or $K = \langle x^d \rangle$ where $d$ is smallest positive integer such that $x^d \in K$. 
    \item If $|H| = \infty$, then for any distinct nonnegative integers $a, b$, $\langle x^a \rangle \neq \langle x^b \rangle$. Futhermore, for every integer $m$, $\langle x^m \rangle = \langle x^{|m|} \rangle$, so that the nontrivial subgroups of $H$ correspond bijectively with the integers $1, 2, 3, \cdots$. 
    \item If $|H| = n < \infty$, then for each positive integer $a$ dividing $n$ there is a unique subgroup of $H$ of order $a$. This subgroup is cyclic group $\langle x^d \rangle$ where $d = n / a$. Futhermore, for every integer $m$, $\langle x^m \rangle = \langle x^{(n, m)} \rangle$ so that the subgroups of $H$ correspond bijectively with the positive divisors of $n$. 
\end{enumerate}

\end{tcolorbox}

\begin{proof}
    
\begin{enumerate}
    \item For such $x^d \in K$, $\langle x^d \rangle \leq K$ is obvious. However, if there exists $x^m \in K$ such that $x^m \notin \langle x^d \rangle$, 
    then we can find $k \in \BZ$ such that $0 < m - dk < d$, and since $x^{m - dk} \in K$, this is contradiction to choice of $d$. 
    \item Easy to show. 
    \item Suppose that $K$ be any subgroup of order $a$. Then by (1), $K = \langle x^b \rangle$. Then here, 
    \[\frac{n}{d} = a = |K| = |x^b| = \frac{n}{(n, b)}\]
    This implies that $d = (n, b)$. So $d | b$ and implies that $x^b \in \langle x^d \rangle$. Hence $\langle x^d \rangle \leq \langle x^b \rangle$. However they have 
    same order, $K = \langle x^d \rangle$. 
\end{enumerate}

\end{proof}

\begin{tcolorbox}[colback=yellow!10!white,colframe=violet!75!white,title=Exercise Sec 2.3.,breakable]
\end{tcolorbox}

\chapter{Quotient Groups, Homomorphisms}

\section{Definitions}

The ultimate goal of group theory is understanding the structure of groups, and then classify them. For that, 
we need to understand \lq smaller' groups contained in original group. We discuss about subgroups in previous section. 
Actually, there are another meaningful smaller groups, namely \textit{quotient groups}.  

\begin{tcolorbox}[colback=yellow!10!white,colframe=brown!75!black,title=Fibers and Homomorphism]

Suppose that $\varphi: G \rightarrow H$ is a group homomorphism. The \textit{fibers} of $\varphi$ are 
the sets of elements of $G$ projecting to single elements of $H$. Actually, $H$ can be intuitively \lq larger' group than $G$. 
However, the image of $G$, namely $\varphi(G)$ is at most large as $G$. So, a \textit{fiber} is
elements of $G$ such that zipped into one points of $\varphi(G)$. 
For any $a \in \varphi(G) \subseteq H$, we can define $X_a \subseteq G, \forall x \in X_a, \varphi(x) = a$ as a \textit{fiber above} $a$.  
\end{tcolorbox}

We'll dive into these interesting subgroups of $G$ during this chapter. 

\begin{tcolorbox}[colback=yellow!10!white,colframe=brown!75!black,title=Kernel]

Let $\varphi : G \rightarrow H$ be a homomorphism. The \textit{kernel} of $\varphi$ is 
\[\ker \varphi := \{g \in G ~|~ \varphi(g) = 1_H\}\]
i.e., \textit{kernel} is fiber above $1_H$. 
\end{tcolorbox}

\begin{tcolorbox}[colback=yellow!10!white,colframe=green!75!black,title=Proposition 1]

Let $G, H$ be groups and $\varphi : G \rightarrow H$ be a homomorphism. 

\begin{enumerate}
    \item $\varphi(1_G) = 1_H$
    \item $\varphi(g^{-1}) = \varphi(g)^{-1}$, for all $g \in G$
    \item $\varphi(g^n) = \varphi(g)^n$, for all $g \in G$
    \item $\ker \varphi \leq G$ 
    \item $\varphi(G) \leq H$
\end{enumerate}

\end{tcolorbox}

\begin{tcolorbox}[colback=yellow!10!white,colframe=brown!75!black,title=Quotient Groups]

Let $\varphi : G \rightarrow H$ be a homomorphism with $\ker \varphi = K$. The \textit{quotient group}, $G / K$ is the group 
whose elements are the fibers of $\varphi$ with the group operation : $X_a \cdot X_b = X_{ab}$

\end{tcolorbox}

In this definition, we intuitively know how they are working. But one is difficult to formalize these things and 
use as a tool to understand the group structure of $G$. So, we'll view quotient group and homomorphisms as alternating perspective. 

\begin{tcolorbox}[colback=yellow!10!white,colframe=green!75!black,title=Proposition 2]

Let $\varphi: G \rightarrow H$ be a homomorphism and $\ker \varphi = K$. For a fiber on $a$, $\varphi^{-1}(a), X_a \in G / K$, 
\[\forall u \in X_a, X_a = \{uk ~|~ k \in K\} = \{ku ~|~ k \in K\}\]

\end{tcolorbox}

Means that, we can enumerate all elements of any fibers by \textbf{one element} and \textbf{kernel}. 

\begin{tcolorbox}[colback=yellow!10!white,colframe=brown!75!black,title=Cosets]

For any $N \leq G$ and $g \in G$, 
\[gN = \{g\cdot n ~|~ n \in N\}, \quad Ng = \{n \cdot g ~|~ n \in N\}\]
called respectively a \textit{left coset} and \textit{right coset} of $N$ in $G$. 
Any element of a coset is called a \textit{representative} for the coset. 
\end{tcolorbox}

\begin{prop}
    
For homomorphism $\varphi: G \rightarrow H$ with $K = \ker \varphi$, 
every fiber is a left(right) coset of $K$ in $G$. And, if $N \leq G$ is kernel of some homomorphism, $gN = Ng$ for any $g \in G$. 
To represent coset, the choice of element does not matter. Means that, $gN = g'N$ if $g, g'$ are in same coset. 

\end{prop}

This proposition gives us that, when we discuss quotient groups, it can be viewed as a set of left cosets. 

\begin{tcolorbox}[colback=yellow!10!white,colframe=blue!75!black,title=Theorem 3]

Let $G$ be a group and $K$ is the kernel of some(unknown) homomorphism from $G$. Then, 
collection of every left cosets of $K$ in $G$: 
\[\{gK ~|~ \forall g \in G\}\]
with following binary operation: 
\[uK \cdot vK = (uv)K\]
forms group, $G / K$. Sometimes we denote $\bar{u}$ to coset $uK$, the coset of $K$ containing $u$.

\end{tcolorbox}

This theorem is quite nice since we can understand $G/K$ without explicit homomorphism and computation of fibers \textbf{if $K$ is kernel of some homomorphism}. 
Then natural question is that, what can we talk about \textbf{kernel of some homomorphism}? 

\begin{tcolorbox}[colback=yellow!10!white,colframe=green!75!black,title=Proposition 4]

Let $N \leq G$. The set of left cosets of $N$ in $G$ forms a partition of $G$. Furthermore, $\forall u, v \in G, uN = vN$ if and only if $v^{-1}u \in N$ if and only if $u, v$ are in the same coset.

\end{tcolorbox}

\begin{proof}
    
For any $g \in G$, there exists coset $gN \subseteq G$, so 
\[\bigcup_{g \in G} gN = G\]
Now it is enough to show that, all distinct cosets are disjoint. Suppose that $g_1N \neq g_2N$. 
If $\exists x \in g_1N \cap g_2N$, $x = g_1n_1 = g_2n_2$ by definition. Then, $g_1 = g_2 n_2 n_1^{-1} = g_2m$, here $m \in N$ since $N \leq G$
so $g_1 N \subseteq g_2 N$ and vice versa, $g_2 N \subseteq g_1 N$. This is contradiction, so $g_1 N \cap g_2N = \emptyset$. 

\end{proof}

\begin{tcolorbox}[colback=yellow!10!white,colframe=green!75!black,title=Proposition 5]

Let $N \leq G$. 

\begin{enumerate}
    \item The operation on set of left cosets of $N$ in $G$ described by 
    \[uN \cdot vN = (uv)N\]
    is well-defined if and only if $gng^{-1} \in N$ for all $g \in G$ and all $n \in N$. 
    i.e., $N_G(N) = G$. 
    \item If the above operation is well-defined, then it makes group $G/N$. 
\end{enumerate}

\end{tcolorbox}

\begin{proof}
    
($\Rightarrow$) The well-definedness implies that 
\[\forall u, u_1 \in uN, \forall v, v_1 \in vN, \quad uvN = u_1v_1N \]
Then, $n, 1\in N, g^{-1} \in g^{-1}N$, we can obtain from proposition 4: 
\[ng^{-1}N = g^{-1}N \iff gng^{-1} \in N\]

($\Leftarrow$) Suppose that $N_G(N) = G$. Then, $\forall u, u_1 \in uN, u_1 = un$ and similarly, $v_1 = vm$. 
Then, 
\[u_1 v_1 = (un)(vm) = u(vv^{-1})nvm = uv(v^{-1}nv)m = uvm' \in uvN\]

\end{proof}

\begin{tcolorbox}[colback=yellow!10!white,colframe=brown!75!black,title=Normal Subgroup]

For $N \leq G$, if $N_G(N) = \{g \in G ~|~ gng^{-1} \in N\} = G$ then $N$ is \textit{normal subgroup} of $G$ and denote 
$N \trianglelefteq G$

\end{tcolorbox}

\begin{prop}
    $G / N$ is subgroup of $G$ if and only if $N \trianglelefteq G$ if and only if $gN = Ng$ for all $g \in G$. 
\end{prop}

\begin{prop}
    To determine that $N \trianglelefteq G$, if we know $N = \langle A \rangle$, it is enough to show that 
    $gag^{-1} \in N$ for any $a \in A, g \in G$. Further, if we know $G = \langle B \rangle$, it is sufficient to show that 
    $bnb^{-1} \in N$ for any $n \in N, b \in B$. 
\end{prop}

\begin{tcolorbox}[colback=yellow!10!white,colframe=green!75!black,title=Proposition 7]

$N \trianglelefteq G \iff N = \ker \varphi \iff G / N$ is group.

\end{tcolorbox}

\begin{defn}

Here, we can easily find above $\varphi: G \rightarrow G/N$ by $g \mapsto gN$. This is obviously 
homomorphism whose kernel is $N$, is called \textit{natural homomorphism} of $G$ onto $G/N$. 

\end{defn}

\begin{tcolorbox}[colback=yellow!10!white,colframe=red!75!black,title=!Caution!]

Note that $A \trianglelefteq B$, $B \trianglelefteq C$ does not implies that $A \trianglelefteq C$. 
i.e., $\trianglelefteq$ does not transitive. The \textit{normality} is relative. 

\end{tcolorbox}

\begin{tcolorbox}[colback=yellow!10!white,colframe=violet!75!white,title=Exercise Sec 3.1.]
\end{tcolorbox}

\newpage 

\section{Cosets and Lagrange's Theorem}

In this section, we'll analyze the subgroup structure of $G$. One of the most important feature of group is its order. 
One important theorem is following: 

\begin{tcolorbox}[colback=yellow!10!white,colframe=blue!75!black,title=Lagrange's Theorem]

$G$ is finite group and $H \leq G$. Then, $|H|$ divide $|G|$ and the number of left cosets of $H$ in $G$ is $\frac{|G|}{|H|}$.

\end{tcolorbox}

\begin{proof}
    
It is enough to show that the order of every coset is same. Then, since $1H = H$ also a coset and cosets form a partition of $G$, 
claim holds. Then we'll find bijection: 
\[\forall g \in G, \quad H \rightarrow gH\]
This is injective, since if $gh_1 = gh_2$ then $h_1 = h_2$ by cancellation in $G$. 
Also surjective, obvious from the definition of $gH$. So clearly this is bijection, so 
\[\forall g \in G, \quad |gH| = |H|\]
This completes the proof.

\end{proof}

\begin{tcolorbox}[colback=yellow!10!white,colframe=brown!75!black,title=Index]

Let $G$ be a group (possibily infinite) and $H \leq G$, the number of left cosets of $H$ in $G$ 
is called \textit{index} of $H$ in $G$, and denoted by $|G:H|$

\end{tcolorbox}

\begin{tcolorbox}[colback=yellow!10!white,colframe=green!75!black,title=Corollary 9]

Let $G$ be a finite group and $x \in G$, then $|x|$ divides $|G|$, so $x^{|G|} = 1$

\end{tcolorbox}

\begin{proof}
    
Since $x \in G$, we can construct $H = \langle x \rangle$ and $H \leq G$. Since $G$ is finite, 
$|H|$ divide $|G|$. Since $H$ is cyclic group, $|H| = |x|$, so $|x|$ divides $|G|$. 

\end{proof}

\begin{tcolorbox}[colback=yellow!10!white,colframe=green!75!black,title=Corollary 10]

    Let $G$ be a finite group and $|G| = p$ is prime, then $G \cong Z_p$
    
\end{tcolorbox}

\begin{proof}
    
By above, for $x \in G$, $|x| = 1$ or $p$ since $p$ is prime. If $|x| = |G| = p$, then 
\[G = \{1, x, \cdots, x^{p-1}\}\]
Then $G$ is cyclic group ($=\langle x \rangle$), so $G \cong Z_p$. 

\end{proof}

\begin{tcolorbox}[colback=yellow!10!white,colframe=gray!75!black,title=Example,breakable]

Let $G$ be any group and $H \leq G$, $|G:H| = 2$. Then, $H \trianglelefteq G$

\end{tcolorbox}

\begin{proof}
    
Since $|G:H| = 2$, there are 2 cosets of $H$ exists: $H$ and $G - H = gH$. 
Suppose that there exists $h \in H$ such that $ghg^{-1} \in gH$. Then it implies that 
\begin{align*}
    ghg^{-1} = gh' & \implies hg^{-1} = h' \in H \\
    & \implies g = h'^{-1}h \in H
\end{align*}
However, $g \in gH = G - H$. This is contradiction, so $\forall g \in G - H$ and $\forall h \in H$, $ghg^{-1} \in H$. 
For $g \in H$, it is obvious. So, $N_G(H) = G$ and $H \trianglelefteq G$. 

\end{proof}

\begin{tcolorbox}[colback=yellow!10!white,colframe=red!75!black,title=Remarks]

\begin{itemize}
    \item \lq Normal' does not means common. Normal subgroup is quite rare structure when $G$ is given. 
    \item Moreover, finding normal subgroup is hard problem. 
    \item If $G$ is finite group and $n$ divides $|G|$, then $G$ need not have a subgroup of order $n$.
    \item If $G$ is abelian and $n$ divides $|G|$, then subgroup whose order is $n$ must exists. 
\end{itemize}

\end{tcolorbox}

\begin{tcolorbox}[colback=yellow!10!white,colframe=blue!75!black,title=Cauchy's Theorem]

If $G$ is a finite group and $p$ is a prime dividing $|G|$ then $G$ has an element of order $p$. 

\end{tcolorbox}

\begin{proof}
    


\end{proof}

\begin{tcolorbox}[colback=yellow!10!white,colframe=blue!75!black,title=Sylow's Theorem]

    If $G$ is a finite group of order $p^\alpha$, where $p$ is a prime and $p$ does not divide $m$, then $G$ has a subgroup of order $p^\alpha$.
    
\end{tcolorbox}

Now, we introduce another very natural construction. 

\begin{tcolorbox}[colback=yellow!10!white,colframe=brown!75!black,title=Subgroup Multiplication]

Let $H, K \leq G$. We define 
\[HK = \{hk ~|~ h \in H, k \in K\}\]

\end{tcolorbox}

\begin{tcolorbox}[colback=yellow!10!white,colframe=green!75!black,title=Proposition 13]

If $H, K$ are finite subgroup then 
\[|HK| = \frac{|H||K|}{|H \cap K|}\]

\end{tcolorbox}

\begin{proof}
    
We can denote 
\[HK = \bigcup_{h \in H} hK\]

We know that each cosets have same order $|K|$. So it is enough to show that there are 
$\frac{|H|}{|H\cap K|}$ distinct cosets. We already show that 
\begin{align*}
    h_1K = h_2K &\iff h_2^{-1}h_1 \in K \\
    & \iff h_2^{-1}h_1 \in K \cap H \\
    & \iff h_1(H \cap K) = h_2(H \cap K)
\end{align*}

Then, since $H, K \leq G$ implies that $H \cap K \leq H$, so number of distinct cosets forms $hK$ is 
same with $|H:H\cap K| = \frac{|H|}{|H\cap K|}$. 

\end{proof}

\begin{tcolorbox}[colback=yellow!10!white,colframe=red!75!black,title=!Caution!]

We can't guarantee $HK$ be a subgroup in general. See following proposition.

\end{tcolorbox}

\begin{prop}
    
$H, K \leq G$, $HK \leq G \iff HK = KH$

\end{prop}

\begin{proof}
    
($\Rightarrow$) If $HK \leq G$, then for any $k \in K, h \in H$, obtain $k, h \in HK$ so $kh \in HK$. 
So, $KH \subseteq HK$. Conversely, let $hk \in HK$. Then $k^{-1}h^{-1} = h'k' \in HK$. 
\[hk = (h'k')^{-1} = k'^{-1}h'^{-1} \in KH\]
So, $HK \subseteq KH$.

($\Leftarrow$) If $HK = KH$, then for any $hk, h'k' \in HK$, 
\begin{align*}
    (hk)(h'k')^{-1} &= h(kk'^{-1})h'^{-1} \\
    & = h (k_0 h_0) \\
    &= h (h_1 k_1) \\
    & = (h h_1) k_1 \in HK
\end{align*}

\end{proof}

\begin{tcolorbox}[colback=yellow!10!white,colframe=green!75!black,title=Corollary 15(For Second Isomorphism Theorem)]

If $H, K \leq G$ and $H \leq N_G(K)$ then $HK \leq G$. Further, if $K \trianglelefteq G$ then $HK \leq G$ for any $H \leq G$. 

\end{tcolorbox}

\begin{proof}
    
Since $H \leq N_G(K)$, we know that $hKh^{-1} = K$ for any $h \in H$. 
Then, this directly implies that $hK = Kh$, so $HK = KH$. So, $HK \leq G$. If $K \trianglelefteq G$, then for any $H \leq G$, $H \leq N_G(K)$ obtained. This completes the proof.

\end{proof}

\begin{tcolorbox}[colback=yellow!10!white,colframe=violet!75!white,title=Exercise Sec 3.2.]
\end{tcolorbox}

\newpage 

\section{The Isomorphism Theorems}

\begin{tcolorbox}[colback=yellow!10!white,colframe=blue!75!black,title=First Isomorphism Theorem]

If $\varphi: G \rightarrow H$ is a homomorphism, then $\ker \varphi \trianglelefteq G$ and $G/\ker \varphi \cong \varphi(G)$
\[
\begin{tikzcd}[row sep=0.5em]
        G \arrow{rr}{\varphi: \text{ Hom}} \arrow[two heads]{dd} & & H \\
        & \mathcolor{red}{\circlearrowleft} & \\
        G/\ker \varphi \arrow[rr, leftrightarrow, red, "\cong", "\exists! \bar{\varphi}"'] & & \varphi(G) \arrow[hook]{uu}
\end{tikzcd}
\]
\end{tcolorbox}

\begin{tcolorbox}[colback=yellow!10!white,colframe=green!75!black,title=Corollary 17]

Let $\varphi: G \rightarrow H$ be a homomorphism. 
\begin{enumerate}
    \item $\varphi$ is injective if and only if $\ker \varphi = 1$
    \item $|G : \ker \varphi| = |\varphi(G)|$
\end{enumerate}

\end{tcolorbox}

\begin{prop}
$\varphi: V \rightarrow W$ is a linear transformation, then $\dim V = \rank ~ \varphi + \mathrm{nullity} ~ \varphi$
\end{prop}

\begin{proof}
    
\[\varphi(V) \cong V / \ker \varphi \implies \rank ~ \varphi = \dim V - \mathrm{nullity } \varphi\]

\end{proof}

\begin{tcolorbox}[colback=yellow!10!white,colframe=blue!75!black,title=Second Isomorphism Theorem]

If $A, B \leq G$ and $A \subseteq N_G(B)$ then $AB \leq G$ and $B \trianglelefteq AB$, $A\cap B \trianglelefteq A$, and 
\[(AB / B) \cong (A / (A \cap B))\]

\[
\begin{tikzcd}[row sep = 0.5em]
    A \arrow[rr] \arrow[dd, two heads] & & AB \arrow[dd, two heads] \\
    & \mathcolor{red}{\circlearrowleft} & \\
    A/(A \cap B) \arrow[rr, red, leftrightarrow, "\cong", "\exists!"'] & & AB/B
\end{tikzcd}
\]
\[
\begin{tikzcd}[row sep = 0.5em]
    & G \arrow[d, no head, dotted] & \\
    & AB \arrow[ld] \arrow[rd] & \\
    A \arrow[rd] &  & B \arrow[ld] \\
    & A \cap B & \\
    & 1 \arrow[u, no head, dotted] & 
\end{tikzcd}
\]

\end{tcolorbox}

\begin{proof}
    
As corollary 15 mentioned, we can obtain $AB \leq G$ directly. Next, since $A \subseteq N_G(B)$ and $B \subseteq N_G(B)$ is obvious, we can obtain that $AB \subseteq N_G(B)$ and so $B \trianglelefteq AB$. This implies that $AB/B$ is group and $a_1B \cdot a_2B = (a_1 a_2)B$ is well-defined. Then now, let's imagine following mapping: 
\[\varphi: A \rightarrow AB/B\]
work as 
\[\varphi: a \mapsto (a\cdot 1)B = aB\]
Here, $\varphi$ is clearly a homomorphism on $A$. The kernel of this map is 
\[\ker \varphi = \{a \in A ~|~ aB = B\} = \{a \in A ~|~ a \in B\} = A \cap B\]
And $\varphi$ is surjective, so $\varphi(A) = AB/B$. 
So we can apply First Isomorphism Theorem: 
\[
\begin{tikzcd}
    A \arrow[rr, "\varphi", two heads] \arrow[dd, two heads] & & AB/B \\
      & &      \\
    A/(A\cap B) \arrow[rr, red, "\cong", leftrightarrow]& & AB/B \arrow[uu, hook]
\end{tikzcd}
\]

So, 
\[AB/B \cong A/(A \cap B)\]

\end{proof}

\begin{tcolorbox}[colback=yellow!10!white,colframe=gray!75!black,title=Remark]

More often, instead of $A \subseteq N_G(B)$ or $A \leq N_G(B)$, following condition also works. 
\[B \trianglelefteq G\]
(Since this directly implies that $A \leq N_G(B)$)
\end{tcolorbox}

\begin{tcolorbox}[colback=yellow!10!white,colframe=blue!75!black,title=Third Isomorphism Theorem]

Let $G$ be a group and $H \trianglelefteq G$, $K \trianglelefteq G$, $H \leq K$. Then $K/H \trianglelefteq G/H$ and 
\[(G/H)/(K/H) \cong G/K\]

\[
\begin{tikzcd}
    G \arrow[dd, two heads] \arrow[rr, two heads] & & G/H \arrow[dd, two heads] \\
    & \mathcolor{red}{\circlearrowleft} & \\
    G/K \arrow[rr, no head, red, "\cong"]& & (G/H)/(K/H)
\end{tikzcd}
\]

\end{tcolorbox}

\begin{proof}
    
For any $kH \in K/H$ and $gH \in G/H$, since $kH \in G/H$ and $H \trianglelefteq G$, following is well-defined. 
\[(gH)(kH)(gH)^{-1} = (gH)(kH)(g^{-1}H) = (gkg^{-1})H \in K/H\]
So, 
\[K/H \trianglelefteq G/H\]
Then, we can imagine following mapping: 
\[\varphi: G/H \rightarrow G/K\]
works by 
\[\varphi: gH \mapsto gK\]
Then $\varphi$ is homomorphism and 
\begin{align*}
    \ker \varphi &= \{gH \in G/H ~|~ gK = K\} \\
    &= \{gH \in G/H ~|~ g \in K\} \\
    &= K/H 
\end{align*}

So, we can draw First Isomorphism Theorem here: 
\[
\begin{tikzcd}
    G/H \arrow[rr, "\varphi"] \arrow[dd, two heads] & & G/K \\
        & &     \\
    (G/H)(K/H) \arrow[rr, red, leftrightarrow, "\cong"] & & G/K \arrow[uu, hook]
\end{tikzcd}
\]
This shows that 
\[(G/H)/(K/H) \cong G/K\]
\end{proof}

Now, we want to use $G/N$ as tool of understanding structure of $G$. Then one natural consequence is about their lattices. 
We'll observe following in Fourth Isomorphism Theorem.

\[
\begin{tikzcd}
    G \arrow[d, no head, dotted, "\text{Lattice of }G/N"] \\
    N \arrow[d, no head, dotted, "\text{Lattice of }N"']\\
    1
\end{tikzcd}
\]

\begin{tcolorbox}[colback=yellow!10!white,colframe=blue!75!black,title=Fourth Isomorphism Theorem]

Let $N \trianglelefteq G$ then there exists a bijection $A \rightarrow \bar{A}$ between 
\[\{A \leq G ~|~ N \subseteq A\} \leftrightarrow \{\bar{A} = A/N \leq G/N\}\]
And for $N \leq A, B \leq G$, 
\begin{enumerate}
    \item $A \leq B$ if and only if $\overbar{A} \leq \overbar{B}$
    \item $A \leq B$ then $|B : A| = |\overbar{B} : \overbar{A}|$
    \item $\overbar{\langle A, B \rangle} = \langle \overbar{A}, \overbar{B} \rangle$
    \item $\overbar{A \cap B} = \overbar{A} \cap \overbar{B}$
    \item $A \trianglelefteq G$ if and only if $\overbar{A} \trianglelefteq \overbar{G}$
\end{enumerate}
\end{tcolorbox}

\begin{proof}
    
For any $\overbar{A} \leq G/N$, $\overbar{A}$ is collection of cosets and coset binary operation is well-defined. 
So, 
\[\forall a_1N , a_2N \in \overbar{A}, \quad (a_1a_2)N \in \overbar{A}\] 
So, the complete preimage of $\overbar{A}$ is disjoint union of each complete preimage of $aN \in \overbar{A}$, 
noted by 
\[A = \biguplus_{aN \in \overbar{A}} A_a\]
Then, $\forall a, b \in A$, we know that $bN \in \overbar{A}$, so $b^{-1}N \in \overbar{A}$ and so $b^{-1} \in A$. 
So, we can obtain 
\[ab^{-1} \in A \implies A \leq G\]
So, above bijection is surjective and $\overbar{A} = A/N$. 

Now suppose that $\overbar{A_1} = \overbar{A_2}$. Then as we've seen, $A_1/N = A_2/N$. 
It is not hard to see that $A_1 = A_2$. 

\end{proof}

\begin{tcolorbox}[colback=yellow!10!white,colframe=gray!75!black,title=Remarks]

Let $\varphi: G \rightarrow G/N$ be natural homomorphism. Then, for $H \leq G$, 
\[\varphi(H) = \varphi(HN)\]

\end{tcolorbox}

\begin{tcolorbox}[colback=yellow!10!white,colframe=violet!75!white,title=Exercise Sec 3.3.]


    
\end{tcolorbox}

\newpage 

\section{}

\chapter{Group Actions}

\section{}

\section{Group Actions: Left Multiplication}

\section{Group Actions: Conjugation}

\section{Automorphisms}

\section{Sylow's Theorem}

\chapter{Direct and Semidirect Products}

\section{Direct Products}



\end{document}